%! TEX root = main.tex

\section{Perfectly Matched Layer (PML)}
\label{sc:15}

In order to reduce spurious reflections at the right boundary, an alternative to the absorbing (impedance) boundary condition at \(x = 1\) is considered. The computational domain is extended by adding an absorbing layer, referred to as a Perfectly Matched Layer (PML), in the region \(x \in (1, 1 + L_{\mathrm{PML}})\). The PML is designed to damp outgoing waves without reflecting them back into the physical domain \(\Omega = (0,1)\).

\subsection{Damping Profile}

The attenuation mechanism is introduced through a spatially varying damping function \(\sigma(x)\). In order to achieve effective absorption while minimizing numerical reflections due to abrupt coefficient variations, \(\sigma(x)\) is required to increase smoothly from zero at the interface \(x = 1\). A parabolic profile is adopted, defined as

\begin{equation}
    \sigma(x) =
    \begin{cases}
        0 & \quad \text{for } x \leq 1 \\
        \sigma_{\max}\left(\dfrac{x - 1}{L_{\mathrm{PML}}}\right)^{2} & \quad \text{for } x > 1
    \end{cases}
    \label{eq:sigma_pml}
\end{equation}

\subsection{Complex Stretching and Modified Coefficients}

The PML formulation is obtained by introducing the complex stretching function

\begin{equation}
    s(x) = 1 - i\frac{\sigma(x)}{\omega}
    \label{eq:stretching}
\end{equation}

Within the layer, the Helmholtz model is recast by replacing the physical parameters \(\rho\) and \(\mu\) with complex-valued equivalents,

\begin{equation}
    \tilde{\rho}(x) = \rho(x)s(x),
    \qquad
    \tilde{\mu}(x) = \frac{\mu(x)}{s(x)}
    \label{eq:pml_coeffs}
\end{equation}

This choice preserves the structure of the governing equation while introducing attenuation through complex coefficients.

\subsection{Impedance Matching and Exponential Decay}

A key property of the PML is that it is perfectly matched at the interface \(x = 1\). In particular, the characteristic impedance in the PML region satisfies

\begin{equation}
    Z_{\mathrm{PML}}
    = \sqrt{\tilde{\rho}(x)\tilde{\mu}(x)}
    = \sqrt{\rho(x)s(x)\frac{\mu(x)}{s(x)}}
    = \sqrt{\rho(x)\mu(x)}
    = Z
    \label{eq:pml_impedance}
\end{equation}

which shows that the impedance is continuous across the interface. This condition implies, at the continuous level, perfect transmission and therefore the absence of reflected waves at \(x = 1\). On the other hand, the local wavenumber becomes complex. Defining \(k = \omega\sqrt{\rho/\mu}\) as the physical wavenumber, the PML wavenumber is obtained as

\begin{equation}
    k_{\mathrm{PML}}(x)
    = \omega\sqrt{\frac{\tilde{\rho}(x)}{\tilde{\mu}(x)}}
    = \omega\sqrt{\frac{\rho(x)s(x)}{\mu(x)/s(x)}}
    = \omega\sqrt{\frac{\rho(x)}{\mu(x)}}s(x)
    = ks(x)
    \label{eq:pml_wavenumber}
\end{equation}

Since \(s(x) = 1 - i\sigma(x)/\omega\) has a negative imaginary part for \(\sigma(x) \geq 0\), the transmitted wave in the layer experiences an exponential decay in amplitude as \(x\) increases, leading to strong attenuation before reaching the truncated boundary.

\subsection{Boundary Conditions on the Extended Domain}

The excitation at the left boundary is kept unchanged and the Neumann condition is imposed at \(x = 0\). The right boundary is moved to the end of the PML, located at \(x = 1 + L_{\mathrm{PML}}\), where a homogeneous Dirichlet condition is prescribed, \(u(1 + L_{\mathrm{PML}}) = 0\). This truncation is justified by the assumption that the outgoing wave has been sufficiently damped by the PML and is therefore negligible at the outer boundary.

\subsection{Numerical Results}

The PML approach is validated by comparing its solution against the impedance boundary condition result obtained in section \ref{sc:13}, for the same physical setup: \(\mu(x)\) as in equation \eqref{eq:mu_jump}, \(\rho = 1\), \(g_N = 1\) and \(\omega = 5\pi\). The computational domain is extended to \((0,\ 1.2)\) with \(L_{\mathrm{PML}} = 0.2\) and \(\sigma_{\max} = 100\), discretized uniformly with a total of \(N = 240\) elements. Figure \ref{fig:pml_comp} reports the amplitude \(|u(x)|\) and the unwrapped phase \(\arg(u(x))\) over the entire extended domain \((0,\ 1.2)\), allowing the behavior inside the absorbing layer to be observed directly alongside the reference impedance BC solution restricted to \((0,1)\).

\begin{figure}[!htb]
    \centering
    \subfloat[Amplitude \(|u(x)|\). In the physical domain \((0,1)\) the two curves are virtually indistinguishable, confirming the absence of spurious reflections at the interface \(x = 1\). In the absorbing layer \((1,\ 1.2)\) the PML solution exhibits rapid exponential decay toward zero, demonstrating the effectiveness of the complex damping coefficient \eqref{eq:pml_coeffs}]{\includegraphics[width=0.90\textwidth]{00_Images/HW1/5.1_amplitude.png}}
    \\
    \subfloat[Unwrapped phase \(\arg(u(x))\). In the physical domain \((0,1)\) the phase profiles of the impedance BC and PML solutions coincide exactly. Inside the absorbing layer \((1,\ 1.2)\) the phase of the PML solution continues to decrease linearly, a signature of an outgoing traveling wave whose amplitude is simultaneously being attenuated by the complex stretching function \eqref{eq:stretching}]{\includegraphics[width=0.90\textwidth]{00_Images/HW1/5.2_phase.png}}
    \caption{Amplitude \(|u(x)|\) and unwrapped phase \(\arg(u(x))\) for the impedance BC solution (Point 3, section \ref{sc:13}) and the PML solution (Point 5, section \ref{sc:15}) over the extended domain \((0,\ 1.2)\). The dashed vertical line marks the physical--PML interface at \(x = 1\)}
    \label{fig:pml_comp}
\end{figure}

\subsection{Physical Interpretation}

The results in figure \ref{fig:pml_comp} demonstrate the two distinct roles played by the PML: it acts as a transparent boundary for the physical domain while simultaneously providing strong attenuation of the transmitted field in the layer region.

\paragraph{Physical domain \(0 \leq x \leq 1\): perfect agreement} The amplitude and phase of the PML solution overlap exactly with those of the impedance BC reference throughout the entire physical domain. The oscillatory interference pattern for \(x < 0.5\) and the smooth transmitted-wave profile for \(x > 0.5\) are reproduced identically by both approaches. This confirms that the impedance-matching property \eqref{eq:pml_impedance} is effectively realized at the discrete level, and that no spurious reflection is introduced at the interface \(x = 1\).

\paragraph{Absorbing layer \(1 < x \leq 1.2\): exponential decay} Within the PML, the amplitude decays exponentially and reaches values negligibly close to zero well before the outer boundary at \(x = 1.2\), validating the effectiveness of the parabolic damping profile \eqref{eq:sigma_pml} with \(\sigma_{\max} = 100\). This rapid attenuation justifies the application of the homogeneous Dirichlet condition at the truncation point without introducing significant error into the physical domain.

\paragraph{Absorbing layer \(1 < x \leq 1.2\): linear phase decay} While the amplitude vanishes, the unwrapped phase of the PML solution continues to evolve linearly and monotonically inside the layer. This behavior is physically crucial: it demonstrates that the wave inside the PML retains the character of an outgoing traveling wave, propagating in the positive \(x\) direction with the same phase velocity as in the physical domain, while its energy content is simultaneously dissipated by the imaginary part of the complex stretching function \(s(x)\). The coexistence of phase propagation and amplitude attenuation is precisely the defining property of the PML, which distinguishes it from simpler damping strategies that also distort the phase and can therefore introduce reflections.